\section{Background}
\label{sec:background}

We briefly review speculative decoding (Section~\ref{sec:background_specdec}) and related work (Section~\ref{sec:related_work}) to provide important background and context for our work.

\subsection{Speculative Decoding}
\label{sec:background_specdec}

Speculative decoding accelerates sampling from a target distribution \(p_{\mathrm{target}}\) by using a cheaper draft distribution \(p_{\mathrm{draft}}\) to generate candidate continuations, and then using the target distribution to accept selectively.
To verify, the target first computes probabilities for all the speculated tokens in parallel in \textit{a single forward pass}.
Drafted tokens are accepted sequentially, each with probability
\(\min\{1,\,p_{\mathrm{target}}(x)/p_{\mathrm{draft}}(x)\}\).
Upon the first rejection, remaining draft tokens are discarded, and the last-token target logits are used to sample a ``bonus token.''
This is done by sampling the bonus token from a modified distribution that guarantees the resulting sequence is distributed as \(p_{\mathrm{target}}\).
% This is done by modifying the last-token target logits in a specific way that guarantees the resulting sequence is distributed as \(p_{\mathrm{target}}\).
This modified distribution is called the \textbf{\textit{residual distribution}}, and takes the form \(r(\cdot) \propto \max\left(p_{\mathrm{target}}(\cdot)-p_{\mathrm{draft}}(\cdot), 0\right)\).
In the case where all tokens are accepted, the residual distribution is simply the target distribution.

The expected number of generated tokens per round, and thus the overall efficiency of speculative decoding, is governed by the \textbf{\textit{acceptance rate}}:
the probability of accepting a given token in the speculation, conditioned on accepting all prior tokens.
In \citep{leviathan2023}, it is shown that the acceptance rate can be written in terms of how well the draft distribution approximates the target distribution in the following way. 

\begin{theorem} \citep{leviathan2023}
    \label{thm:l1}
    $$\alpha = \sum_x \min\{p_{\mathrm{target}}(x),\,p_{\mathrm{draft}}(x)\} = 1 - \tfrac{1}{2}\lVert p_{\mathrm{target}} - p_{\mathrm{draft}} \rVert_1$$
\end{theorem}

\begin{definition}
We define a \textbf{speculation} at round $T$, denoted $\spec^T \coloneqq (s_1^T, \ldots, s_K^T)$, as the sequence of $K$ tokens proposed autoregressively by the draft model.
The length $K$ of the speculated sequence is called the \textbf{speculative lookahead}.
\end{definition}

\begin{definition}
We define a \textbf{verification outcome} at round $T$, denoted $\ver^T \coloneqq (k, t^*) \in \Ver^T$, where $(s_1^{T-1}, \ldots, s_k^{T-1})$ are the accepted draft tokens from round $T-1$, and $t^*$ is the \textbf{bonus token} sampled from the residual distribution (or target distribution if all tokens are accepted). 
\end{definition}

% \color{red}
\subsection{Related Work}
\label{sec:related_work}
% We now discuss prior work that, like SSD, performs speculation and verification in parallel to one another, as well as work that allow verifying full trees of tokens.

\textbf{Parallel Speculative Decoding.}
AMUSD~\citep{amusd} and PEARL~\citep{pearl} propose speculating the next round during ongoing verification, but only prepare for the verification outcome in which all tokens are accepted, which is one special case of many possible outcomes.
SwiftSpec~\citep{swiftspec} and SpecBranch~\citep{specbranch} prepare a larger cache consisting of a token tree branching off of the speculation being verified (thus enabling larger speedups), but both use fallback strategies that do not work at large batch sizes. SpecBranch~\citep{specbranch} has motivation similar to SSD's: it constitutes (approximately) a special case of the SSD framework where only a single branching point is allowed, where the fallback speculator is equal to the regular speculator, and where the hyperparameters (branching point, number branches, speculation length) are dynamically chosen to maximize speedups.
SwiftSpec~\citep{swiftspec} considers the special case of greedy sampling, and adopts a fallback strategy of just-in-time speculation that struggles at higher temperatures and batch sizes when cache-misses become inevitable.

\textbf{Tree-Based Speculative Decoding.}
Numerous speculative decoding methods have been proposed that increase the expected number of accepted tokens by allowing the draft model to speculate a \textit{tree} of tokens instead of a sequence, thereby giving the verifier several token options at each position~\citep{specinfer,eagle2,sequoia,specexec}.
Our method differs in important ways:
First, existing tree-based methods are still sequential: speculate, then verify.
Second, tree-based methods introduce a large amount of \textit{verifier compute}---which is quite expensive due to the size of the target model---because the entire tree must be verified in the target model forward pass.
Our method, on the other hand, scales up the \textit{speculation compute} by pre-speculating for many verification outcomes in parallel, but does not introduce any additional verifier compute.
Lastly, it is important to note that our method can be combined with these tree-based approaches for further gains (see Appendix~\ref{app:combine} for discussion). 

\textbf{Improved Draft Architectures.}
There have been many advancements in draft model architectures that can improve the acceptance rates and/or speeds of the draft model and which can be used in SSD for even better performance.
For example, EAGLE~\citep{eagle,eagle2,eagle3} allows the draft model to take as input the powerful representations from the target model for the current prefix, while GliDe~\citep{du2024glidecapelowhasslemethod} and LongSpec~\citep{yang2025longspeclongcontextlosslessspeculative} allow the draft model to perform cross-attention on the KV cache of the target model.
Alternate model architectures, like diffusion LLMs~\citep{Nie2025LargeLanguageDiffusionModels, sahoo2024simple,diffuspec,specdiff,llmfuture,dflash} or SSMs~\citep{gu2022efficiently, Gu2024Mamba,mamballama}, can also be used to increase the speed with which the draft model can produce the speculated token sequence.
We include more technical details on how SSD can fruitfully be combined with these improved draft model architectures (e.g., EAGLE-3; \cite{eagle3}) in Appendix \ref{app:combine}.

\color{black}